%!TEX root = ../thesis.tex
\chapter{Methodology}
\label{chap:methodology}

%% Approx. 15 pages.
%% Main thrusts: exploratory lightweight-head SSL, cross-city SSL benchmark,
%% DiffusionDrive + SSL, and mechanistic analysis.
%% Ends with the Phase 0--4 experimental protocol.

This chapter is organised to mirror the narrative of the thesis. It starts with the lightweight-head agent, because that is the simplest setting in which representation quality can be isolated from planner capacity. It then explains why Drive-JEPA changed the scope of the project from pooled NAVSIM performance to city-disjoint transfer, and why TransFuser and Latent TransFuser are the main closed-loop anchor planners for the cross-city benchmark. The chapter next introduces DiffusionDrive as a second planner family with stronger built-in priors, and it closes with the mechanistic methods used later to explain the observed SSL gains.

% ============================================================================
\section{Exploratory Lightweight-Head Regime}
\label{sec:meth-lightweight-head}

The project began with an intentionally simple experimental regime: pair a pretrained visual backbone with a lightweight planning head and ask whether the representation alone already improves end-to-end driving behavior. This regime is deliberately minimal. By holding planner capacity fixed at its lowest meaningful level and varying only the backbone, any change in downstream performance is attributable to the representation rather than to the planner. Understanding the lightweight regime is a prerequisite for the cross-city experiments in~\ref{sec:meth-city-splits} and the generative-head study in~\ref{sec:meth-diffusiondrive}.

\subsection{Role in the Thesis}
\label{sec:meth-lh-role}

Three questions motivated the exploratory regime. First, whether the in-distribution improvement of SSL over supervised backbones reported in related work~\cite{wang-drive-jepa, naeinian-cross-city-2026} survives when the planner is stripped to a waypoint regressor. Second, whether different SSL objectives (I-JEPA, DINOv2, MAE) separate from one another under the same head, or whether their differences require a heavier planner to become visible. Third, whether fine-tuning the backbone helps or hurts relative to freezing it, once the planning head is this small. The answers to those questions set the prior for every subsequent planner choice in the thesis.

The answer from Drive-JEPA~\cite{wang-drive-jepa} reframed the scope. Their result that V-JEPA plus a lightweight decoder is already competitive on standard NAVSIM argued that in-distribution PDMS on the pooled benchmark is close to saturation for SSL methods. The exploratory regime therefore serves less as a competitive submission and more as a controlled baseline: it produces the head-fixed, backbone-varying numbers against which the cross-city and diffusion-head studies measure their own contributions.

\subsection{Agent Architecture}
\label{sec:meth-lh-arch}

The exploratory planner is a single configurable agent, \texttt{LightweightPlanningAgent}, whose components are wired through a dataclass configuration so that a backbone, a multi-view fusion rule, and a planning head can be swapped independently. \reffig{fig:method-lh-arch} shows the forward path.

\begin{figure}[ht]
    \centering
    \resizebox{\linewidth}{!}{%
    \begin{tikzpicture}[
        node distance = 4mm and 5mm,
        box/.style = {draw, rounded corners=2pt, minimum height=8mm, minimum width=18mm, align=center, font=\scriptsize},
        ssl/.style = {fill=black!5},
        arrow/.style = {-{Stealth[length=1.8mm]}, thick},
    ]
        \node[box] (cams) {Cameras\\$V \in \{1,3,6\}$};
        \node[box, right=of cams, ssl] (enc) {Backbone\\(ResNet / ViT)};
        \node[box, right=of enc] (tok) {Tokens\\$B{\times}V{\times}N{\times}C$};
        \node[box, right=of tok, ssl] (fus) {View fusion\\$f_\text{view}$};
        \node[box, right=of fus] (vis) {Visual\\$B{\times}C$};
        \node[box, below=8mm of fus] (ego) {Ego\\$B{\times}8$};
        \node[box, right=of vis, ssl] (head) {Head\\(MLP / Tx.)};
        \node[box, right=of head] (traj) {Trajectory\\$B{\times}T{\times}3$};
        \draw[arrow] (cams) -- (enc);
        \draw[arrow] (enc) -- (tok);
        \draw[arrow] (tok) -- (fus);
        \draw[arrow] (fus) -- (vis);
        \draw[arrow] (vis) -- (head);
        \draw[arrow] (ego) -| (head);
        \draw[arrow] (head) -- (traj);
    \end{tikzpicture}}
    \caption{The \texttt{LightweightPlanningAgent} forward path. A camera set is encoded by a shared backbone into per-view patch tokens, fused into a single visual vector by one of five interchangeable rules, concatenated with ego status, and decoded to $T{=}8$ waypoints. Shaded blocks are the components swapped during ablations; unshaded blocks are fixed.}
    \label{fig:method-lh-arch}
\end{figure}

Every block in \reffig{fig:method-lh-arch} is addressable by configuration. The camera set is selected by a \texttt{camera\_mode} field that resolves to one, three, or six nuPlan views. The backbone is constructed by the same \texttt{create\_image\_encoder} factory used in the TransFuser and DiffusionDrive variants of~\ref{sec:meth-backbones}; it accepts any backbone in \reftab{tab:backbones} and exposes an \texttt{unfreeze\_last\_blocks} hook for encoder-fine-tune ablations. The view-fusion module $f_\text{view}$ consumes a $(B, V, N, C)$ token tensor and returns a $(B, C)$ summary; five implementations are available. The head consumes $[\text{visual} \;\|\; \text{ego}]$ and emits eight poses. Keeping the interface narrow at each seam lets the experiment matrix vary one factor at a time.

\subsection{Design Decisions}
\label{sec:meth-lh-design}

Four design axes are load-bearing and are reused later in the chapter. Each axis is exposed as a named field on \texttt{LightweightHeadConfig} and selected by a single Hydra override. \reftab{tab:meth-lh-design} enumerates the choice space and the rationale for each axis.

\begin{table}[ht]
\centering
\caption{Configuration axes of the lightweight-head agent. Each field is a top-level key in \texttt{LightweightHeadConfig}; values are the Hydra strings used at the command line.}
\label{tab:meth-lh-design}
\small
\renewcommand{\arraystretch}{1.2}
\setlength{\tabcolsep}{4pt}
\begin{tabular}{@{}>{\raggedright\arraybackslash}p{3.2cm} >{\raggedright\arraybackslash}p{3.3cm} >{\raggedright\arraybackslash}p{1.7cm} >{\raggedright\arraybackslash}p{4.3cm}@{}}
\toprule
\textbf{Axis} & \textbf{Values} & \textbf{Default} & \textbf{Rationale} \\
\midrule
\texttt{backbone} &
  \texttt{resnet}, \texttt{ijepa}, \texttt{dinov2}, \texttt{mae} &
  \texttt{ijepa} &
  Single \texttt{create\_} \texttt{image\_encoder} factory shared with Trans\-Fuser and Diffusion\-Drive. \\
\texttt{vit\_backend} &
  \texttt{huggingface}, \texttt{custom} &
  \texttt{hugging\-face} &
  \texttt{custom} loads Meta-style checkpoints with exact 2-D sin--cos positional encodings. \\
\texttt{trainable\_}\newline\texttt{backbone\_}\newline\texttt{fraction} &
  $\{0.0,\ 1.0\}$ &
  $0.0$ &
  Frozen vs.\ fully-unfrozen only; intermediate fractions caused optimizer instability in pilots. \\
\texttt{camera\_mode} &
  \texttt{front1}, \texttt{front3}, \texttt{front5\_}\texttt{plus\_b} &
  \texttt{front1} &
  Selects 1, 3, or 6 nuPlan cameras. Heavier sets reserved for fusion ablations. \\
\texttt{view\_fusion} &
  \texttt{single}, \texttt{mean}, \texttt{weighted\_}\texttt{mean}, \texttt{concat\_}\texttt{proj}, \texttt{cross\_attn} &
  \texttt{single} &
  \texttt{cross\_attn} adds ${\approx}\,21$\,M parameters and restricts the GPU pool. \\
\texttt{head\_type} &
  \texttt{mlp}, \texttt{transformer} &
  \texttt{mlp} &
  \texttt{mlp} is the two-layer $1288{\to}256{\to}256{\to}T{\cdot}3$ regressor used for representation isolation. \\
\bottomrule
\end{tabular}
\end{table}

Two decisions that cut across the axes in \reftab{tab:meth-lh-design} are worth calling out explicitly. First, when \texttt{trainable\_backbone\_fraction}${=}\,1.0$, the optimizer places encoder parameters in a separate group with a lower learning rate than the head, matching the two-group AdamW schedule used in~\ref{sec:meth-phase1-law}. Second, the fusion rules are defined on the same $(B, V, N, C)\to(B, C)$ signature, which means the fusion choice is a clean independent axis in the experiment matrix rather than a backbone-specific option. Both decisions carry forward into the heavier planners later in this chapter.

\subsection{Exploratory Experiment Taxonomy}
\label{sec:meth-lh-taxonomy}

The exploratory regime emits a structured experiment matrix rather than a list of hand-tuned runs. Four tiers organize the matrix; they run in priority order so that the main thesis claim is answered before any expensive variant is scheduled. \reffig{fig:method-lh-matrix} summarizes the tiers.

\begin{figure}[ht]
    \centering
    \begin{tikzpicture}[
        node distance = 7mm and 10mm,
        tier/.style = {draw, rounded corners=3pt, minimum height=10mm, minimum width=80mm, align=center, font=\small},
        t1/.style = {fill=black!10},
        t2/.style = {fill=black!5},
        arrow/.style = {-{Stealth[length=1.8mm]}, semithick, gray!70, shorten >=7pt, shorten <=3pt},
    ]
        \node[tier, t1] (p1) {\textbf{Priority 1}\\11 backbones $\times$ frozen, MLP, front1 (10 runs)};
        \node[tier, t2, below=of p1] (p2) {\textbf{Priority 2}\\4 backbones $\times$ 4 training cities + 4 fusion-ablation variants (20 runs)};
        \node[tier, below=of p2] (p3) {\textbf{Priority 3}\\11 backbones $\times$ fully-unfrozen, MLP, front1 (10 runs)};
        \node[tier, below=of p3] (p4) {\textbf{Priority 4}\\3-cam / 6-cam / cross-attn / transformer-head, fully-unfrozen (5 runs)};
        \draw[arrow] (p1) -- (p2);
        \draw[arrow] (p2) -- (p3);
        \draw[arrow] (p3) -- (p4);
    \end{tikzpicture}
    \caption[Four-tier exploratory experiment matrix (45 runs).]{Four-tier exploratory experiment matrix, totalling $45$ runs. Tier order is enforced: the Priority~1 backbone-controlled baseline completes before heavier variants are scheduled, so that a fully-unfrozen ViT-H/14 cannot overwrite the backbone-only comparison. Priority~2 extends across the four training cities and adds fusion ablations on a fixed I-JEPA backbone; Priority~3 repeats Priority~1 with the encoder fully fine-tuned; Priority~4 reserves the transformer-head and multi-camera variants. Heavier tiers are pinned to higher-memory GPUs because fully-unfrozen ViT-H/14 backbones do not fit on 48\,GB cards.}
    \label{fig:method-lh-matrix}
\end{figure}

Priority~1 is the core scientific instrument of this section. Each backbone in \reftab{tab:backbones} is paired with the same frozen-encoder \texttt{mlp} head and evaluated on the pooled \texttt{navtest} split, so the only variable is backbone identity. Priority~2 extends the protocol across the four nuPlan cities introduced in~\ref{sec:meth-city-splits} and adds fusion ablations on a fixed I-JEPA backbone. Priority~3 repeats the Priority~1 matrix with the encoder fully unfrozen, while Priority~4 reserves the heavier transformer and multi-camera variants. Across all runs, the head, loss, optimizer schedule, camera normalization, and trajectory sampling are held fixed (see \reftab{tab:meth-lh-recipe}).

\subsection{Training and Evaluation Recipe}
\label{sec:meth-lh-recipe}

All 45 runs share the single recipe in \reftab{tab:meth-lh-recipe}. The recipe is intentionally minimal: no auxiliary losses (world-model, semantic, detection) are used in this regime, because any non-zero auxiliary weight would confound the representational comparison. The only hyperparameters that vary between runs are the encoder learning rate (which takes effect only when the encoder is unfrozen) and the batch size (which drops to 16 for fully-unfrozen or multi-camera runs because of memory).

\begin{table}[ht]
\centering
\caption{Fixed training and evaluation recipe for the exploratory regime.}
\label{tab:meth-lh-recipe}
\small
\renewcommand{\arraystretch}{1.15}
\setlength{\tabcolsep}{6pt}
\begin{tabular}{@{}>{\raggedright\arraybackslash}p{3.8cm} >{\raggedright\arraybackslash}p{8.5cm}@{}}
\toprule
\textbf{Setting} & \textbf{Value} \\
\midrule
Input resolution      & $224 \times 224$, normalized with per-backbone processor statistics \\
Cameras               & 1, 3, or 6 (from \texttt{camera\_mode}) \\
Trajectory target     & 8 future poses at 0.5\,s (4\,s horizon), ego-frame $(x, y, \text{heading})$ \\
Loss                  & $L_1$ trajectory regression; no auxiliary losses \\
Optimizer             & AdamW, weight decay $0.0$ \\
Head LR               & $1 \times 10^{-4}$ \\
Encoder LR            & $3 \times 10^{-5}$ (applied only when the encoder is unfrozen) \\
LR schedule           & Cosine, decays to $1\%$ of initial over 20 epochs \\
Epochs                & 20 \\
Precision             & Mixed (bfloat16 via PyTorch Lightning) \\
Gradient clipping     & $1.0$ \\
Batch size            & 32 (frozen, single-camera) / 16 (unfrozen or multi-camera) \\
Host memory per run   & $128$\,GB (batch $32$) / $256$\,GB (batch $16$) \\
Eval split            & \texttt{navtest} \\
Eval scorer           & one-stage PDM (non-reactive) \\
Devkit                & NAVSIM v1.0 primary; v1.1 re-evaluation \\
Primary metric        & PDMS (sub-scores NC, DAC, EP, TTC, Comfort also logged) \\
\bottomrule
\end{tabular}
\end{table}

The split, scorer, and devkit choice match the protocol in~\cite{naeinian-cross-city-2026}, which keeps the numbers in Chapter~\ref{chap:experiments} directly comparable to that prior work. The \texttt{v1\_1} re-eval is a defensive control for the eval-version sensitivity discussed in Chapter~\ref{chap:relatedwork}.

\subsection{From Lightweight Head to Cross-City}
\label{sec:meth-lh-bridge}

Three observations carry forward from this regime. The lightweight \texttt{mlp} head is sufficient to discriminate backbones: with planner capacity held at its minimum, backbone choice alone moves in-distribution PDMS by several points, as reported in~\ref{sec:exp-phase0}. This means the cross-city study in the following sections can use a similarly minimal planner without risking that the reported gaps are masked by planner expressivity. The fusion interface $(B, V, N, C) \to (B, C)$ generalizes to the heavier planners later in the chapter, so the same multi-camera variants can be revisited without modifying the planning head. Finally, the same shared configuration schema drives all $45$ runs; extending the matrix with four training cities in~\ref{sec:meth-city-splits} is a matter of populating additional rows, not changing the protocol.

The remaining open question, and the one~\ref{sec:meth-city-splits} takes up, is whether the backbone ordering produced by this regime on a pooled benchmark survives geographic distribution shift, or whether SSL's apparent advantage collapses once train and test cities differ. That question cannot be answered at the level of \texttt{navtest}; it requires a re-partitioning of the benchmark along city identity, which the next section builds.

% ============================================================================
\section{NAVSIM City Splits}
\label{sec:meth-city-splits}

\subsection{City Extraction from nuPlan Metadata}
\label{sec:meth-city-extraction}

City labels are recovered directly from the underlying nuPlan/OpenScene log metadata. Each scenario token inherits the city prefix of its source log, with prefixes such as \texttt{us-ma-boston}, \texttt{us-nv-las-vegas}, \texttt{us-pa-pittsburgh}, and \texttt{sg-one-north}. A preprocessing script walks the metadata once, assigns every scenario token to exactly one city, and writes city-specific token lists used by the training and evaluation dataloaders. This step is simple but essential: it turns NAVSIM from a pooled benchmark into a reproducible geographic-transfer benchmark whose train/test partitions are defined by city identity rather than random sampling.

\subsection{Split Statistics and Distribution}
\label{sec:meth-split-stats}

The cross-city protocol uses the natural geographic partitions already present in the source datasets. \reftab{tab:meth-split-stats} gives the exact split sizes; the prose interpretation is simple: NAVSIM is substantially imbalanced, and Singapore is both the smallest training source and the only left-hand-traffic environment. That combination makes Singapore a stringent transfer case and a confound that must be kept visible when interpreting cross-city results~\cite{naeinian-cross-city-2026}. \reffig{fig:method-city-map} illustrates the four cities schematically; the full data card is in \refapp{app:city-data-card}.

\begin{table}[ht]
\centering
\caption{City-split statistics used by the cross-city protocols.}
\label{tab:meth-split-stats}
\small
\renewcommand{\arraystretch}{1.15}
\setlength{\tabcolsep}{6pt}
\begin{tabular}{@{}llrrr@{}}
\toprule
\textbf{Benchmark} & \textbf{City} & \textbf{Train scenes} & \textbf{Eval scenes} & \textbf{Train hours} \\
\midrule
nuScenes & Boston & 390 & 77 & 2.12 \\
nuScenes & Singapore & 310 & 73 & 1.68 \\
\midrule
NAVSIM & Las Vegas & 818 & 50 & -- \\
NAVSIM & Boston & 255 & 62 & -- \\
NAVSIM & Pittsburgh & 140 & 22 & -- \\
NAVSIM & Singapore & 97 & 13 & -- \\
\bottomrule
\end{tabular}
\end{table}

\begin{figure}[ht]
    \centering
    \begin{tikzpicture}[
        city/.style = {draw, rounded corners=3pt, minimum width=26mm, minimum height=14mm, align=center, font=\small},
        continent/.style = {draw, densely dashed, rounded corners=4pt, inner sep=4mm, font=\scriptsize\itshape, label position=above, label distance=-1mm},
        citycount/.style = {font=\scriptsize, gray},
    ]
        \node[city, fill=red!8] (lv) {Las Vegas (LV)\\{\scriptsize\color{gray}818 / 50}};
        \node[city, fill=red!8, right=10mm of lv] (pit) {Pittsburgh (P)\\{\scriptsize\color{gray}140 / 22}};
        \node[city, fill=red!8, right=10mm of pit] (bos) {Boston (B)\\{\scriptsize\color{gray}255 / 62}};
        \node[draw, densely dashed, rounded corners=4pt, inner sep=3mm, fit=(lv) (pit) (bos), label=above:\itshape North America, label distance=-1mm] (na) {};

        \node[city, fill=blue!8, right=18mm of bos] (sg) {Singapore (S)\\{\scriptsize\color{gray}97 / 13}\\\tiny left-hand traffic};
        \node[draw, densely dashed, rounded corners=4pt, inner sep=3mm, fit=(sg), label=above:\itshape Southeast Asia, label distance=-1mm] {};

        \node[font=\scriptsize, below=4mm of lv.south west, anchor=north west, text width=120mm] {\textit{trainval / test scene counts from \refapp{app:city-data-card}. Singapore is the only left-hand-traffic city and the smallest training source, making it a stringent cross-city target.}};
    \end{tikzpicture}
    \caption{Schematic of the four NAVSIM City Splits. Three North-American cities share traffic conventions (right-hand-side driving) and broadly comparable road-marking standards; Singapore adds a different traffic convention and distinct road-paint conventions. Scene counts are from \reftab{tab:city-splits}.}
    \label{fig:method-city-map}
\end{figure}

\subsection{Evaluation Protocols}
\label{sec:meth-eval-protocols}

The thesis uses two evaluation settings. The all-city protocol trains and evaluates on the benchmark's original mixed-city split and anchors numbers to published baselines. The stress test is strict single-city transfer: a model is trained on one city and evaluated on unseen cities with no fine-tuning, adaptation, or test-time modification. On nuScenes this yields a two-direction protocol (Boston$\rightarrow$Singapore and Singapore$\rightarrow$Boston); on NAVSIM it yields a full $4 \times 4$ train-city by test-city matrix. Leave-one-city-out evaluation is noted as a future extension. The core empirical results come from the first two settings, matching~\cite{naeinian-cross-city-2026}.
\begin{enumerate}
    \item \textbf{All-city (standard)}: Train and test on the full NAVSIM split (all cities mixed). Reproduces published numbers.
    \item \textbf{Single-city (zero-shot transfer)}: Train on city A, test on all cities. Produces a $4 \times 4$ generalization matrix.
    \item \textbf{Leave-one-city-out (LOCO)}: Train on 3 cities, test on the held-out city. More realistic multi-city regime.
\end{enumerate}

For score-based metrics such as PDMS, we summarize out-of-distribution performance for a training city $c_i$ by averaging over the other cities~\cite{naeinian-cross-city-2026}:
\begin{equation}
    S_\text{OOD}(c_i)=\frac{1}{|C|-1}\sum_{c_j \neq c_i} S(c_i \rightarrow c_j).
\end{equation}

We additionally report the generalization gap as the difference between in-domain and out-of-domain closed-loop performance:
\begin{equation}
    \Delta_\text{gen} = \text{PDMS}_\text{in-dist} - \text{PDMS}_\text{OOD}.
\end{equation}

For open-loop error metrics, where lower is better, the more informative quantity is the error inflation ratio
\begin{equation}
    R_\text{err}(c_\text{train}\rightarrow c_\text{test})=\frac{E_\text{cross}(c_\text{train}\rightarrow c_\text{test})}{E_\text{in}(c_\text{train}\rightarrow c_\text{train})},
\end{equation}
which equals 1 for perfect preservation and grows as transfer degrades~\cite{naeinian-cross-city-2026}.

% ============================================================================
\section{SSL Pre-Training Protocol}
\label{sec:meth-ssl-pretraining}

All domain-specific SSL pretraining used in this thesis was carried out by Fatemeh Naeinian, first author of our earlier cross-city benchmark study~\cite{naeinian-cross-city-2026}, on a shared pretraining repository maintained for this collaboration. This section documents the recipe she followed, which is the protocol under which every nuScenes-pretrained ViT-S/14 backbone in Chapters~\ref{chap:methodology} and~\ref{chap:experiments} was produced. The goal of the protocol is not to maximize absolute benchmark performance during pretraining; it is to create matched representation initializations whose downstream behavior can be compared fairly under city transfer. Holding the architecture, data, schedule, and input resolution fixed across three SSL families is what turns the backbone comparisons in Chapter~\ref{chap:experiments} into a controlled intervention rather than an uncontrolled family comparison.

\subsection{Common Setup: Dataset, Architecture, and Input Geometry}
\label{sec:meth-ssl-common}

The three pretraining recipes share a common substrate. The pretraining source is the nuScenes training split~\cite{caesar-nuscenes-2020}, expanded to all six camera rigs (\texttt{CAM\_FRONT}, \texttt{CAM\_FRONT\_LEFT}, \texttt{CAM\_FRONT\_RIGHT}, \texttt{CAM\_BACK}, \texttt{CAM\_BACK\_LEFT}, \texttt{CAM\_BACK\_RIGHT}) rather than the front camera alone. This yields \(168{,}780\) RGB frames. Using every camera matters for two reasons. It increases the pretraining signal roughly six-fold relative to a front-only protocol, and it exposes the encoder to the full driving field of view that downstream planners will see, including side and rear cameras whose statistics differ from front-facing ImageNet pretraining. All frames are normalized with the ImageNet RGB mean and standard deviation; no colour jitter, Gaussian blur, or heavy photometric augmentation is applied, because such augmentations bias the invariance set of the representation toward classification tasks and are not load-bearing for driving planning.

The model architecture is ViT-S/14~\cite{dosovitskiy-vit-2021}. Its concrete size, data source, resolution choices, and hardware envelope are summarized in \reftab{tab:meth-ssl-common}. The design choice is methodological rather than leaderboard-oriented: the domain-specific runs intentionally use a smaller architecture than the large public ImageNet checkpoints, so that the experiments can ask whether driving-domain pretraining compensates for scale under geographic transfer.

Input geometry is treated as a controlled factor rather than as an implementation detail. The square crop is the ImageNet-compatible baseline; the rectangular crop preserves more horizontal driving context. Training both variants to completion lets the transfer study separate objective effects from geometry effects.

\begin{table}[ht]
\centering
\caption{Shared pretraining substrate for all three SSL objectives (I-JEPA, DINOv2, MAE).}
\label{tab:meth-ssl-common}
\small
\renewcommand{\arraystretch}{1.2}
\setlength{\tabcolsep}{6pt}
\begin{tabular}{@{}>{\raggedright\arraybackslash}p{3.8cm} >{\raggedright\arraybackslash}p{8.5cm}@{}}
\toprule
\textbf{Setting} & \textbf{Value} \\
\midrule
Source dataset            & nuScenes~\cite{caesar-nuscenes-2020} training split \\
Camera views              & All six (\texttt{CAM\_FRONT}, \texttt{CAM\_FRONT\_\{LEFT,RIGHT\}}, \texttt{CAM\_BACK}, \texttt{CAM\_BACK\_\{LEFT,RIGHT\}}) \\
Frames                    & 168{,}780 RGB images \\
Normalization             & ImageNet mean / std \\
Photometric augmentation  & None (no color jitter, no Gaussian blur) \\
Backbone                  & ViT-S/14 (patch 14, $D{=}384$, $L{=}12$, $6$ heads) \\
Input resolutions         & $224 \times 224$ (square) and $224 \times 392$ (rectangular) \\
Epochs                    & 100 \\
Optimizer                 & AdamW with cosine schedule \\
Hardware                  & 1$\times$ H200 (80\,GB) per run \\
\bottomrule
\end{tabular}
\end{table}

\subsection{Method-Level Comparison}
\label{sec:meth-ssl-compare}

The three objectives share the substrate above and differ only in the pretext task and its surrounding machinery. I-JEPA predicts target-block latent representations from visible context, with an EMA target encoder and a predictor; downstream consumers load the EMA target encoder. DINOv2 uses student--teacher distillation with multi-crop augmentation and an EMA teacher; downstream transfer uses the teacher backbone. Patch-level iBOT loss is disabled because its benefit does not transfer cleanly to the small driving corpus. MAE has no EMA teacher; the encoder is extracted after training on masked-patch pixel reconstruction with a lightweight decoder. \reftab{tab:meth-ssl-methods} summarises the per-method settings along the axes that matter for transfer.

\begin{table}[ht]
\centering
\caption{Domain-specific SSL pretraining configurations. All three runs use ViT-S/14 on nuScenes (168{,}780 frames, six cameras) for 100 epochs on one H200 GPU.}
\label{tab:meth-ssl-methods}
\small
\renewcommand{\arraystretch}{1.2}
\setlength{\tabcolsep}{4pt}
\begin{tabular}{@{}>{\raggedright\arraybackslash}p{2.4cm} >{\raggedright\arraybackslash}p{3.4cm} >{\raggedright\arraybackslash}p{3.1cm} >{\raggedright\arraybackslash}p{3.6cm}@{}}
\toprule
\textbf{Axis} & \textbf{I-JEPA} & \textbf{DINOv2} & \textbf{MAE} \\
\midrule
Objective &
  Predict target-block latents from context &
  Match student-teacher softmax over prototypes &
  Reconstruct masked pixels \\
Loss domain &
  Latent space (MSE) &
  Cross-entropy on prototype distribution &
  Pixel space (MSE) \\
Masking / crops &
  1 context block ($85$--$100\%$); 4 target blocks ($15$--$20\%$) &
  2 global crops + 6 local $96{\times}96$ crops &
  $75\%$ random patch mask \\
Teacher &
  EMA, momentum $0.996\!\to\!1.0$ &
  EMA, base momentum $0.9985$ &
  None \\
Objective-specific head &
  Predictor over target-block latents &
  Prototype head; student / teacher temperatures $0.1$ / $0.04$ &
  Decoder depth $8$, dim $512$, $16$ heads \\
Base LR &
  $1\!\times\!10^{-3}$ &
  $2\!\times\!10^{-3}$ &
  $1.5\!\times\!10^{-4}$ \\
Warmup &
  40 epochs &
  10 epochs &
  40 epochs \\
LR decay / final LR &
  Cosine to $1\!\times\!10^{-6}$ &
  Cosine decay &
  Cosine to $0$ \\
Weight decay &
  $0.04 \to 0.4$ linear &
  $0.04 \to 0.4$ linear &
  $0.05$ \\
Precision / clipping &
  -- &
  FP16; clip $3.0$ &
  -- \\
Batch size &
  128--512 (H200) &
  768 (H200) &
  768 (H200) \\
Checkpoint file &
  \texttt{.pth.tar} &
  \texttt{.pth} &
  \texttt{.pth} \\
Checkpoint key &
  \texttt{target\_encoder} &
  \texttt{teacher} (backbone) &
  \texttt{model} (encoder) \\
\bottomrule
\end{tabular}
\end{table}

\subsection{Rectangular vs.\ Square Inputs}
\label{sec:meth-ssl-rect}

The rectangular-input variants exist because the camera streams used downstream are much wider than square ImageNet crops. \reftab{tab:meth-rect-inputs} records the exact geometry and implementation changes. The prose takeaway is that rectangular pretraining preserves horizontal driving context, while square pretraining remains the compatibility control for ImageNet-style ViT checkpoints and prior work~\cite{wang-drive-jepa}.

\begin{table}[ht]
\centering
\caption{Square versus rectangular ViT input geometry.}
\label{tab:meth-rect-inputs}
\small
\renewcommand{\arraystretch}{1.15}
\setlength{\tabcolsep}{6pt}
\begin{tabular}{@{}>{\raggedright\arraybackslash}p{3.4cm} >{\raggedright\arraybackslash}p{4.1cm} >{\raggedright\arraybackslash}p{4.6cm}@{}}
\toprule
\textbf{Axis} & \textbf{Square control} & \textbf{Rectangular variant} \\
\midrule
Resolution & $224 \times 224$ & $224 \times 392$ \\
Patch grid at $P{=}14$ & $16 \times 16$ & $16 \times 28$ \\
Field-of-view effect & ImageNet-compatible crop; can discard horizontal context & Preserves more lane, sidewalk, and adjacent-traffic context \\
Positional embeddings & Standard square grid & Regenerated at rectangular grid, not warped from square \\
Code path & Standard ViT factory & \texttt{vit\_backend}${=}\,\texttt{custom}$ in the rectangular forks \\
\bottomrule
\end{tabular}
\end{table}

\subsection{Reference ImageNet Checkpoints}
\label{sec:meth-ssl-imagenet}

For the large-scale ImageNet-pretrained baselines (I-JEPA ViT-H/14, DINOv2 ViT-L/14, MAE ViT-L/16) the thesis uses the original authors' released weights without additional pretraining. Including these baselines controls for the confound that any observed SSL advantage on nuScenes might simply reflect model scale rather than pretraining-domain match. Their exact model identifiers and feature dimensions are reused verbatim from \reftab{tab:backbones}.

% ============================================================================
\section{SSL Backbone Integration}
\label{sec:meth-backbones}

\subsection{Backbone Adapter Interface}
\label{sec:meth-backbone-adapter}

All backbones are wrapped behind a common adapter interface so that planner code does not depend on model-specific preprocessing or tensor shapes. Conceptually, every backbone implements a common mapping
\[
\texttt{images} \mapsto \texttt{features},
\]
while the adapter handles input normalization, resizing, token pooling, and feature reshaping. This abstraction is important for the thesis because it makes backbone swapping a controlled intervention rather than a code fork. In practice, the same planning stack can be re-run with ResNet, I-JEPA, DINOv2, or MAE by changing configuration rather than rewriting the downstream model.

\subsection{Backbone Configurations}
\label{sec:meth-backbone-configs}

Table~\ref{tab:backbones} summarizes the backbone families used throughout the thesis. The key control principle is to vary the pretraining objective while keeping downstream planners fixed wherever possible.

\begin{table}[ht]
\centering
\caption{SSL backbone configurations.}
\label{tab:backbones}
\begin{tabular}{lllcl}
\toprule
\textbf{Backbone} & \textbf{SSL Type} & \textbf{Pre-train Data} & \textbf{Output Dim} & \textbf{Used In} \\
\midrule
ResNet-34         & Supervised      & ImageNet-1K & 512  & LAW, TF, DD \\
I-JEPA ViT-H/14  & Predictive      & ImageNet-1K & 1280 & LAW, TF \\
I-JEPA ViT-S/14  & Predictive      & nuScenes    & 384  & DD \\
DINOv2 ViT-L/14  & Discriminative  & LVD-142M    & 1024 & LAW, TF \\
MAE ViT-L/16     & Generative      & ImageNet-1K & 1024 & LAW, TF, DD \\
\bottomrule
\end{tabular}
\end{table}

TF denotes TransFuser and DD denotes DiffusionDrive. Unless stated otherwise, SSL backbones are frozen during downstream training. When fine-tuning is enabled for ablations, the backbone receives a smaller learning-rate multiplier than the planning head to reduce catastrophic drift away from the pretrained representation.

\subsection{Feature Projection}
\label{sec:meth-feature-proj}

Backbones emit features with different dimensions and layouts. A lightweight projection standardises them before they enter the planner. For regression heads this is a linear projection from the backbone output dimension to the planner hidden dimension. For token-based backbones in DiffusionDrive, the adapter reshapes ViT patch tokens $(B, N, C)$ into spatial feature maps $(B, C', H, W)$ compatible with the multi-scale cross-attention blocks. The adapter is deliberately thin so that the observed effects attribute to the representation rather than to task-specific adaptation. \reffig{fig:method-adapter} sketches the data flow.

\begin{figure}[ht]
    \centering
    \resizebox{\linewidth}{!}{%
    \begin{tikzpicture}[
        node distance = 4mm and 5mm,
        box/.style = {draw, rounded corners=2pt, minimum height=7mm, minimum width=20mm, align=center, font=\scriptsize},
        arrow/.style = {-{Stealth[length=1.8mm]}, thick},
    ]
        \node[box] (img) {Image\\$3 \times H_0 \times W_0$};
        \node[box, right=of img, fill=black!5] (vit) {ViT backbone\\(frozen)};
        \node[box, right=of vit] (tok) {Patch tokens\\$B{\times}N{\times}C$};
        \node[box, right=of tok] (rsh) {Reshape +\\linear proj.};
        \node[box, right=of rsh, fill=black!5] (map) {Spatial map\\$B{\times}C'{\times}H{\times}W$};
        \node[box, right=of map] (plan) {Planner\\(TF / LAW / DD)};
        \draw[arrow] (img) -- (vit);
        \draw[arrow] (vit) -- (tok);
        \draw[arrow] (tok) -- (rsh);
        \draw[arrow] (rsh) -- (map);
        \draw[arrow] (map) -- (plan);
    \end{tikzpicture}}
    \caption{Backbone adapter schematic. A frozen ViT encoder emits $N$ patch tokens; a lightweight reshape and linear projection restructure them into a spatial feature map whose shape matches what the downstream planner expects. ResNet-34 backbones already emit a spatial map and therefore skip the reshape step. This adapter is the only module that differs between the TransFuser/LAW cross-city runs and the DiffusionDrive runs.}
    \label{fig:method-adapter}
\end{figure}

% ============================================================================
\section{Zero-Shot Cross-City Transfer}
\label{sec:meth-phase1-cc}

The cross-city benchmark is the controlled representation-versus-geography study introduced in our earlier cross-city work~\cite{naeinian-cross-city-2026}, extended here with the full set of backbone configurations from~\ref{sec:meth-ssl-pretraining} and integrated through the same factory described in~\ref{sec:meth-backbones}. Holding the planner and training recipe fixed, does the choice of visual backbone change how a model behaves when it is trained in one city and evaluated in a different one? The protocol is split across two planner families: an open-loop latent-world-model planner on nuScenes~\cite{caesar-nuscenes-2020} and closed-loop multi-modal or camera-only planners on NAVSIM~\cite{dauner-navsim-2024}. Running both reduces the risk that the finding is contingent on a single planner's inductive bias.

The collaboration split of labour matters for attribution. The open-loop LAW pipeline was implemented and run by Fatemeh Naeinian (first author of~\cite{naeinian-cross-city-2026}); it is described here only to the extent needed to interpret its numbers. The closed-loop NAVSIM pipeline, including the TransFuser / Latent TransFuser backbone-swap infrastructure and the configuration-hardening work that these variants required, is the author's contribution and is treated at more detail.

\subsection{Open-Loop Pipeline: LAW on nuScenes}
\label{sec:meth-phase1-law}

The open-loop leg uses LAW~\cite{li-law-2024}, an external baseline, with its stock Swin-Transformer-Tiny image encoder replaced by each of the pretrained ViTs described in~\ref{sec:meth-backbones}. The planner (per-view spatial decoders, LID depth encoding, delta-waypoint regression, and the auxiliary latent world-model head) is held fixed across backbone variants; the only manipulated factor is the encoder initialisation. This design follows~\cite{naeinian-cross-city-2026}. \reftab{tab:meth-phase1-train} summarises the training recipe.

\subsection{Closed-Loop Pipeline: TransFuser and Latent TransFuser on NAVSIM}
\label{sec:meth-phase1-tf}

The closed-loop leg evaluates the same representational question under NAVSIM's pseudo-simulation scorer, where the predicted trajectory is rolled out against log-replay agents and PDMS is computed from the resulting rollout (Section~\ref{sec:bg-navsim}). Two architectures run side by side. \emph{TransFuser}~\cite{chitta-transfuser-2023} keeps its ResNet-34 LiDAR branch but replaces the ResNet-34 image branch with the SSL-pretrained encoder under test; the cross-modal fusion stack is unchanged. \emph{Latent TransFuser} is the same model with the LiDAR branch replaced by a fixed learned latent~\cite{dauner-navsim-2024}; this is the camera-only probe, and the critical variant for testing whether the representation effect survives without the geometric-sensing shortcut. Running both isolates the contribution of LiDAR from the contribution of the backbone.

Plugging a ViT encoder into a fusion stack that was designed around a CNN feature pyramid is not a one-line change. The image encoder must produce a four-stage pyramid whose channel dimensions match what the fusion stack consumes; the adapter used here is the one described in~\ref{sec:meth-feature-proj}. Position embeddings are regenerated at the native camera-crop grid (per the \texttt{vit\_backend}${=}\,\texttt{custom}$ path) rather than bicubically interpolated from a square pretraining, which matters for the rectangular nuScenes-pretrained ViT-S/14 variants. None of these changes alter the fusion layers, the waypoint decoder, or the training loss.

The single non-trivial methodological note in this leg is a faithfulness issue specific to TransFuser. Several defaults in modern Lightning / Hydra scaffolding silently degrade the original recipe if imported unchanged: scheduler insertion, AdamW decay, mismatched LiDAR-branch depth, and accumulator double-counting all move the baseline enough to invalidate a backbone comparison. For that reason, the NAVSIM v1.1 re-evaluation sidecar preserves the original recipe verbatim; the exact optimizer, batch, epoch, and hardware settings are listed in \reftab{tab:meth-phase1-train}.

\begin{table}[ht]
\centering
\caption{Training recipes for the two cross-city pipelines. The open-loop recipe matches~\cite{li-law-2024,naeinian-cross-city-2026}; the closed-loop recipe matches the original TransFuser / Latent TransFuser setup in~\cite{chitta-transfuser-2023,dauner-navsim-2024}.}
\label{tab:meth-phase1-train}
\small
\renewcommand{\arraystretch}{1.2}
\setlength{\tabcolsep}{5pt}
\begin{tabular}{@{}>{\raggedright\arraybackslash}p{3.3cm} >{\raggedright\arraybackslash}p{4.2cm} >{\raggedright\arraybackslash}p{4.6cm}@{}}
\toprule
\textbf{Setting} & \textbf{LAW (open-loop, nuScenes)} & \textbf{TransFuser / LatTF (closed-loop, NAVSIM)} \\
\midrule
Optimizer              & AdamW                            & Adam \\
Base LR                & $5 \times 10^{-5}$                & $1 \times 10^{-4}$ (constant) \\
Weight decay           & $0.01$                            & $0$ \\
Scheduler              & Cosine, min-LR ratio $10^{-3}$    & None (no warmup) \\
Gradient clip (L2)     & $35$                              & $1.0$ \\
Epochs                 & $12$                              & $20$ (cheap) / $100$ (full) \\
Batch size             & $2$ (single GPU)                  & $32$ per GPU (effective $128$) \\
Hardware               & $1 \times$ NVIDIA L40S            & $4 \times$ NVIDIA H200 \\
Precision              & Mixed (bfloat16)                  & Mixed (bfloat16) \\
Backbone frozen/full   & Both evaluated                    & Both evaluated \\
\bottomrule
\end{tabular}
\end{table}

\subsection{Backbone Matrix}
\label{sec:meth-phase1-matrix}

The backbone matrix is the one built in~\ref{sec:meth-ssl-pretraining} and~\ref{sec:meth-backbone-configs}. \reftab{tab:meth-phase1-backbones} is the primary record of the variants; the important design rule is that each row is evaluated under both frozen and fully-unfrozen adaptation modes, so representation quality and representation adaptability can be separated.

\begin{table}[ht]
\centering
\caption{Backbone matrix used in the cross-city benchmark. Every row is evaluated under both adaptation modes (frozen and fully unfrozen) in each of the two planners, giving a controlled substrate for comparison.}
\label{tab:meth-phase1-backbones}
\small
\renewcommand{\arraystretch}{1.2}
\setlength{\tabcolsep}{5pt}
\begin{tabular}{@{}>{\raggedright\arraybackslash}p{3.0cm} >{\raggedright\arraybackslash}p{2.3cm} >{\raggedright\arraybackslash}p{2.7cm} >{\raggedright\arraybackslash}p{3.9cm}@{}}
\toprule
\textbf{Family} & \textbf{Capacity} & \textbf{Pre-training source} & \textbf{Role in the benchmark} \\
\midrule
Swin-T~\cite{liu-swin-2021}         & $T$        & ImageNet-1k (supervised) & LAW stock baseline \\
ResNet-34~\cite{he-resnet-2016}     & $34$       & ImageNet-1k (supervised) & TransFuser / LatTF stock baseline \\
I-JEPA~\cite{assran-ijepa-2023}     & ViT-H/14   & ImageNet-1k (SSL)         & Large-scale SSL reference \\
DINOv2~\cite{oquab-dinov2-2024}     & ViT-S/14   & LVD-142M (SSL)            & Large-scale SSL reference \\
MAE~\cite{he-mae-2022}              & ViT-B/16   & ImageNet-1k (SSL)         & Large-scale SSL reference \\
\midrule
I-JEPA (nuSc, sq) & ViT-S/14 & nuScenes $224{\times}224$ & Domain-matched, square \\
I-JEPA (nuSc, rect) & ViT-S/14 & nuScenes $224{\times}392$ & Domain-matched, rectangular \\
DINOv2 (nuSc, sq) & ViT-S/14 & nuScenes $224{\times}224$ & Domain-matched, square \\
DINOv2 (nuSc, rect) & ViT-S/14 & nuScenes $224{\times}392$ & Domain-matched, rectangular \\
MAE (nuSc, sq) & ViT-S/14 & nuScenes $224{\times}224$ & Domain-matched, square \\
MAE (nuSc, rect) & ViT-S/14 & nuScenes $224{\times}392$ & Domain-matched, rectangular \\
\bottomrule
\end{tabular}
\end{table}

The capacity asymmetry in the top block of \reftab{tab:meth-phase1-backbones} is intentional and follows the convention of our earlier cross-city study~\cite{naeinian-cross-city-2026}: the three ImageNet SSL entries are used \emph{at the capacity they were released at} rather than artificially shrunk to ViT-S/14, because those released weights are what a practitioner would actually deploy; the like-for-like objective comparison instead happens inside the bottom block, where architecture, data, schedule, and input resolution are held fixed across the three SSL families.

\subsection{Geographic Transfer Protocol}
\label{sec:meth-phase1-protocol}

Both legs use strictly city-disjoint training and evaluation. For each training city, no frame from any other city is seen during training, no adaptation or test-time modification is allowed, and no hyperparameter is re-tuned for the target city. Evaluation uses the same splits, pre-processing, and eval scorer as the corresponding in-domain run; the only change is the set of scenarios. \reftab{tab:meth-phase1-splits} enumerates the exact training-test pairs.

\begin{table}[ht]
\centering
\caption{Cross-city transfer matrix. Every cell represents a full training run (train city) evaluated on every evaluation city listed, under the training recipes in \reftab{tab:meth-phase1-train}.}
\label{tab:meth-phase1-splits}
\small
\renewcommand{\arraystretch}{1.2}
\setlength{\tabcolsep}{5pt}
\begin{tabular}{@{}>{\raggedright\arraybackslash}p{2.7cm} >{\raggedright\arraybackslash}p{3.0cm} >{\raggedright\arraybackslash}p{3.1cm} >{\raggedright\arraybackslash}p{2.8cm}@{}}
\toprule
\textbf{Benchmark} & \textbf{Training city} & \textbf{Evaluation city} & \textbf{Pair count} \\
\midrule
nuScenes (LAW) & \{Boston, Singapore\} & \{Boston, Singapore\} & $2 \times 2 = 4$ (two in-domain + two cross-city) \\
\addlinespace
NAVSIM (TF, LatTF) & \{Boston, Las Vegas, Pittsburgh, Singapore\} & Same 4 & $4 \times 4 = 16$ (four in-domain + twelve cross-city) \\
\bottomrule
\end{tabular}
\end{table}

Singapore plays a distinct role on both benchmarks: it is the only left-hand-traffic environment in either dataset, and on NAVSIM it is also the smallest training partition (Section~\ref{sec:meth-split-stats}). The companion cross-city study~\cite{naeinian-cross-city-2026} identifies right-hand to left-hand transfer as the structurally hardest direction, and Chapter~\ref{chap:experiments} reports the directional asymmetry explicitly.

\subsection{Generalization-Gap Metrics}
\label{sec:meth-phase1-metrics}

The in-domain to cross-city behavior is reduced to two scalar summaries that appear throughout Chapter~\ref{chap:experiments}; both follow~\cite{naeinian-cross-city-2026}.

For open-loop error metrics (L2 displacement and collision rate on nuScenes), where lower is better, the transfer inflation from training city \(c_\text{train}\) to evaluation city \(c_\text{test}\) is reported as
\begin{equation}
\label{eq:meth-rerr}
R_\text{err}(c_\text{train} \to c_\text{test})
=
\frac{E_\text{cross}(c_\text{train} \to c_\text{test})}
     {E_\text{in}(c_\text{train} \to c_\text{train})},
\end{equation}
so that \(R_\text{err}=1\) indicates no degradation under transfer and \(R_\text{err}>1\) indicates degradation. Ratios are reported separately for L2 displacement and for collision rate.

For closed-loop score metrics (PDMS on NAVSIM), where higher is better, we report the mean out-of-distribution PDMS of a training city,
\begin{equation}
\label{eq:meth-sood}
S_\text{OOD}(c_i)
=
\frac{1}{|\mathcal{C}|-1}
\sum_{c_j \neq c_i} S(c_i \to c_j),
\end{equation}
which averages the closed-loop scores across every unseen test city. \(S_\text{OOD}\) reduces the \(4\times 4\) NAVSIM matrix to one scalar per training city per backbone, which is the quantity plotted in the per-city panels of Chapter~\ref{chap:experiments}.

Two practical notes. First, \(S_\text{OOD}\) averages over a heterogeneous set of test cities; it is not symmetric in the sense that training on a hard source and evaluating on easy targets can give the same mean as the reverse, so it is always reported alongside the full matrix rather than in isolation. Second, the ratio \(R_\text{err}\) is sensitive to the denominator: when the in-domain run itself has a degenerate score, the ratio becomes uninformative, and such runs are flagged in the tables of Chapter~\ref{chap:experiments} rather than silently included in ranked comparisons.

% ============================================================================
\section{Generative Planning with DiffusionDrive: a Backbone-Attenuating Test}
\label{sec:meth-diffusiondrive}

The cross-city benchmark establishes that the visual representation dominates planner capacity under geographic transfer, but under two planner families (LAW and TransFuser / Latent TransFuser) whose heads are both regression-based and whose cross-modal fusion is a thin attention stack over a convolutional feature pyramid. A natural follow-up question is whether this finding is specific to planners with low expressive capacity, or whether it also holds under a planner whose own priors do most of the work. DiffusionDrive~\cite{jia-diffusiondrive-2025} is selected for this test because its head sits closer to the second category.

\subsection{Why DiffusionDrive is the Right Stress Test}
\label{sec:meth-dd-why}

DiffusionDrive's planner absorbs backbone differences by construction. Four of its design choices make this explicit:
\begin{enumerate}
    \item \textbf{K-means anchor trajectories.} Denoising is initialized from one of 256 prototypical trajectories obtained by K-means clustering on \texttt{navtrain} expert motions~\cite{jia-diffusiondrive-2025}. The planner therefore \emph{starts from a plausible trajectory} on every scenario, rather than from Gaussian noise. A weak backbone whose features cannot resolve fine-grained scene geometry can still exit the denoiser close to the expert trajectory, because the anchor has already placed the plan in the right basin.
    \item \textbf{Truncated two-step denoising.} At inference, DiffusionDrive takes only two denoising steps. The visual conditioning therefore has two opportunities to perturb the anchor before commitment. This is a design decision for real-time throughput~\cite{jia-diffusiondrive-2025}, but its side effect is that each individual conditioning step contributes less than it would under a 50- or 1{,}000-step denoiser.
    \item \textbf{Grid-sampled BEV cross-attention at waypoint positions.} The cross-attention that injects scene features into the denoising query is evaluated on a regular BEV grid around the current waypoint prediction, not on a dense feature map. This quantizes the planner's access to the backbone features and makes the head robust to differences in token geometry.
    \item \textbf{Time-only FiLM conditioning.} The final conditioning path uses FiLM modulation on the denoising timestep only. Backbone-dependent modulations that were explored during development did not end up in the production path; the head therefore sees the backbone features only through the grid-sampled cross-attention, with no additional FiLM gate that might amplify them.
\end{enumerate}

Each of these choices is known from the DiffusionDrive paper~\cite{jia-diffusiondrive-2025} and is preserved verbatim in this thesis. The composite effect is that DiffusionDrive is an \emph{attenuating substrate for the representation claim}: if a different backbone changes planner behavior at all under this head, that change must overcome the anchor-trajectory prior, the two-step budget, and the quantized cross-attention. A planner whose head was already doing most of the work under DiffusionDrive would manifest as an \emph{absence} of backbone effect on this planner, and that would be the cleanest way for the thesis claim to fail.

The consequence for interpretation is that small absolute PDMS differences under DiffusionDrive are more meaningful than similarly-sized differences under a weaker planner. Specifically, the study reports transfer-ratio inflation (in-distribution PDMS divided by out-of-distribution PDMS) rather than raw absolute scores. A ResNet-34 baseline that drops 16 points to Singapore and an SSL variant that drops 14 points is a directionally consistent robustness signal even though the absolute gap is small, because the DiffusionDrive head is actively attenuating the backbone contribution before the PDMS scorer sees the plan.

\subsection{Integration: Same Adapter as the Cross-City Benchmark}
\label{sec:meth-dd-ssl}

The central repro-fairness commitment of the generative-planning study is that the SSL backbone is plugged into DiffusionDrive through the \emph{same} adapter used in the cross-city benchmark. In code, both pipelines import \texttt{create\_image\_encoder} from the TransFuser backbone module (\ref{sec:meth-backbones}); DiffusionDrive's image-encoder configuration is a drop-in replacement at the same interface. There is no bespoke DiffusionDrive-specific ViT adapter, no resolution-specific branch divergence, and no "feature pyramid up-sampling" patch that is applied here but not in the cross-city benchmark.

This matters because it forecloses the obvious alternative explanation for any observed effect: "maybe DiffusionDrive sees the backbone more richly because its adapter is different." The adapter is held fixed, and the head is the single manipulated variable between the cross-city benchmark and this study. A separate ViT-bottleneck ablation branch explores resolution-matching commits that could, in principle, further narrow the backbone--planner gap; those commits are \emph{not} merged into the thesis's main branch and \emph{not} applied to any run reported here, so reproducibility under the same adapter is preserved.

Integration otherwise follows the minimal recipe:
\begin{itemize}
    \item The ResNet-34 image encoder in the published DiffusionDrive configuration is replaced with the SSL backbone under test (I-JEPA, DINOv2, or MAE; HuggingFace IN-1k or nuScenes-pretrained custom-ViT, per \refsec{sec:meth-backbone-configs}).
    \item The ViT patch-token output is reshaped to a four-stage pyramid compatible with DiffusionDrive's \texttt{MultiScaleCrossAttention} via the adapter in \refsec{sec:meth-feature-proj}. Square (\(224{\times}224\)) and rectangular (\(224{\times}392\)) input geometries are both evaluated.
    \item Both frozen and fully-unfrozen adaptation modes are tested, matching the cross-city benchmark protocol. No partial-thaw runs are included.
\end{itemize}

\subsection{Camera-Only Variant}
\label{sec:meth-dd-camera-only}

A camera-only DiffusionDrive variant is included as a stricter probe. LiDAR-dependent inputs and rasterized BEV features are removed; the LiDAR branch is replaced by a fixed learned latent, analogous to the Latent TransFuser construction in~\refsec{sec:bg-transfuser}. The variant targets two readers: the deployment-cost argument for camera-only autonomy, and the stricter test of whether a richer visual representation can partially compensate for the missing geometric channel.

\subsection{Training Configuration}
\label{sec:meth-dd-training}

All DiffusionDrive runs use the shared configuration in \reftab{tab:meth-dd-training}. The interpretation rule is that rows with shortened schedules are architectural probes, not final leaderboard attempts; Chapter~\ref{chap:experiments} flags those rows when reporting results.

\begin{table}[ht]
\centering
\caption{Shared DiffusionDrive training and inference configuration.}
\label{tab:meth-dd-training}
\small
\renewcommand{\arraystretch}{1.15}
\setlength{\tabcolsep}{6pt}
\begin{tabular}{@{}>{\raggedright\arraybackslash}p{4.0cm} >{\raggedright\arraybackslash}p{8.2cm}@{}}
\toprule
\textbf{Setting} & \textbf{Value} \\
\midrule
Optimizer & AdamW \\
Learning rate / decay & $10^{-4}$ with cosine decay \\
Weight decay & $10^{-4}$ \\
Warmup & 3 epochs \\
Training mode & Distributed data parallel where required \\
Anchor bank & 256 K-means trajectories clustered on \texttt{navtrain} \\
Inference denoising & 2 steps \\
Schedule note & Full schedule where compute permitted; shortened all-city probes are flagged in results \\
\bottomrule
\end{tabular}
\end{table}

\subsection{What Counts as a Positive Result Under DiffusionDrive}
\label{sec:meth-dd-success}

Because DiffusionDrive's head is actively attenuating the backbone contribution, the evaluation criterion for the generative-planning and cross-city generative studies is stricter than a "top-of-leaderboard" criterion and weaker than a "beat the ResNet baseline outright" criterion. Specifically, the thesis commits, in advance of running the city-disjoint matrix, to the following success condition:

\begin{quote}
\emph{If frozen SSL backbones produce a smaller OOD PDMS drop than supervised ResNet-34 in three or more of the four training-city rows under DiffusionDrive, the representation-centric claim of the thesis is judged to survive the planner swap under a backbone-attenuating head. If the signal is inconsistent across training cities, the thesis reports a null result in this study and folds the outcome into the limitations section, rather than re-engineering DiffusionDrive to rescue the effect.}
\end{quote}

\noindent
This pre-registered criterion is what the interpretation of the result tables in \refsec{sec:exp-phase3} anchors on.

% ============================================================================
\section{Study Sequence}
\label{sec:meth-exp-protocol}

The experiments are organised as one exploratory study followed by four load-bearing studies. The exploratory lightweight-head study of~\ref{sec:meth-lightweight-head} motivates the representation-first framing. The cross-city SSL benchmark of~\ref{sec:meth-phase1-cc} produces the backbone-by-city matrix that Chapter~\ref{chap:experiments} reports. The generative-planning study adds DiffusionDrive with SSL backbones on the standard NAVSIM split, using the integration described in~\ref{sec:meth-diffusiondrive}. The cross-city generative study extends that integration to the city-disjoint setting; partial rows are reported in Chapter~\ref{chap:experiments}, but the study is not used as load-bearing evidence because the full matched matrix is incomplete. The mechanistic analysis uses CKA, linear probing, MMD, and feature-covariance effective rank on the trained checkpoints, with no additional training.

% ============================================================================
\section{Metrics}
\label{sec:meth-metrics}

Every closed-loop experiment reports overall PDMS and its five sub-scores (NC, DAC, EP, TTC, Comfort). Cross-city runs additionally report the generalisation gap $\Delta_\text{gen}$ and the mean out-of-distribution PDMS $S_\text{OOD}$ defined in~\ref{sec:meth-phase1-metrics}. Open-loop nuScenes results report L2 displacement, collision rate, and the transfer inflation ratio $R_\text{err}$ from the companion cross-city study~\cite{naeinian-cross-city-2026}. DiffusionDrive rows additionally report inference speed.

% ============================================================================
\section{Implementation Details}
\label{sec:meth-implementation}

All experiments were implemented in PyTorch and PyTorch Lightning. Training used NVIDIA L40S and H200 GPUs with distributed data parallelism where required by model size and batch size. Pre-trained checkpoints for the public SSL backbones were loaded from the released model weights and routed through the common adapter described above; all other weights were initialised following the original reference implementation of each planner. The per-study training recipes and evaluation protocols are summarised in Appendix~\ref{app:hparams}.
